\documentclass[11pt]{article}
\usepackage[margin=1in]{geometry}
\usepackage{amsmath,amssymb,amsthm,booktabs,enumitem,graphicx,xcolor,hyperref,microtype}
\usepackage[T1]{fontenc}
\usepackage{lmodern}
\hypersetup{colorlinks=true,linkcolor=blue!55!black,urlcolor=blue!55!black,citecolor=blue!55!black}
\setlist{itemsep=2pt,topsep=4pt}
\newtheorem{theorem}{Theorem}
\newtheorem{lemma}{Lemma}
\newcommand{\C}{\mathbb C}
\newcommand{\Sym}{\operatorname{Sym}}
\newcommand{\id}{\mathrm{id}}
\newcommand{\sgn}{\operatorname{sgn}}
\title{Permutation-Symmetric Ogata Approximation Fails\\for Three Pauli Observables}
\author{Automated mathematics run \texttt{fa\_banach\_001}, lane 17}
\date{11 August 2026}

\begin{document}
\maketitle

\begin{center}
\fbox{\parbox{0.91\linewidth}{\textbf{Status.}
Candidate counterexample, likely valid, requiring human review.  The result is
a full negative answer to the simultaneous tensor-permutation-symmetry
instantiation of the source question; it does not settle every other symmetry
pattern mentioned there.}}
\end{center}

\begin{abstract}
Herrera asks whether the commuting approximants in Ogata's theorem can
preserve the symmetries of three or more macroscopic observables.  We show
that the most canonical tensor-permutation version is false.  For the three
macroscopic Pauli observables on $N$ qubits, every commuting self-adjoint
triple invariant under the transposition $(1\,2)$ and the long cycle
$(1\,2\,\ldots\,N)$ has maximum coordinate error at least $1/6$, independently
of $N$.  The two symmetries generate $S_N$, so the approximants reduce the
fully symmetric tensor block.  There the observables form the normalized
irreducible spin-$N/2$ triple.  Its Pauli localizer has unbalanced inertia,
whereas the invertible localizer of a commuting triple has balanced inertia.
The argument is elementary and gives the explicit constant.
\end{abstract}

\section{Source question and exact scope}

The source is David Herrera, \emph{A Projection Characterization and Symmetry
Bootstrap for Elements of a von Neumann Algebra that are Nearby Commuting
Elements}, arXiv:2412.20795v1 (2024) \cite{Herrera2024}.  In Example 2.5,
PDF pages 14--15, it defines the tensor-factor permutation symmetries
$S_\sigma(A)=U_\sigma A U_\sigma^{-1}$ on $M_d(\C)^{\otimes N}$ and notes that
the macroscopic observable
\[
 T_N(A)=\frac1N\sum_{r=1}^N
 I_d^{\otimes(r-1)}\otimes A\otimes I_d^{\otimes(N-r)}
\]
is fixed by every $S_\sigma$.  It also notes that fixedness for the
transposition $(1\,2)$ and long cycle $(1\,2\,\ldots\,N)$ already implies
fixedness for all permutations.

The question is transcribed as follows:
\begin{quote}
``A natural question that one can ask is if a symmetry version of Ogata's
theorem holds. Namely, when $H_i$ satisfy certain symmetries and
antisymmetries, is it true that the $H_i$ are nearby commuting matrices that
also satisfy these symmetries. How much symmetry that is preserved for
$m\ge 3$ is still not known.''
\end{quote}

\begin{figure}[ht]
\centering
\includegraphics[width=0.94\linewidth]{figures/open_problem_crop.png}
\caption{The open question in arXiv:2412.20795v1, PDF page 15.}
\end{figure}

Our scope is precise: can the two noncommuting permutation generators be
preserved \emph{simultaneously} for the three macroscopic Pauli observables?
The answer is no.  We do not claim that preserving one fixed $S_\sigma$ at a
time is impossible.

\section{Definitions and proof intuition}

Let
\[
 s_1=\frac12\begin{pmatrix}0&1\\1&0\end{pmatrix},\qquad
 s_2=\frac12\begin{pmatrix}0&-i\\i&0\end{pmatrix},\qquad
 s_3=\frac12\begin{pmatrix}1&0\\0&-1\end{pmatrix}
\]
be the norm-$1/2$ Pauli spin matrices, and set
$H_{i,N}=T_N(s_i)$ on $\mathcal H_N=(\C^2)^{\otimes N}$.
Write $U_\sigma$ for the unitary permuting tensor factors according to
$\sigma\in S_N$.

For self-adjoint $X_1,X_2,X_3$ on a finite-dimensional Hilbert space, define
their Pauli localizer
\[
 \mathcal L(X)=\sum_{i=1}^3 X_i\otimes\tau_i,
 \qquad \tau_i=2s_i.
\]
The $\tau_i$ are the unit-norm Pauli matrices.

\paragraph{Proof intuition.}
Ogata's theorem permits its commuting approximants to mix the spin sectors in
the tensor power.  Full permutation invariance removes that escape route.
Indeed, any permutation-invariant operator reduces the symmetric tensor
subspace $K_N=\Sym^N(\C^2)$.  The compressed macroscopic Pauli triple on
$K_N$ is the ``fuzzy sphere,'' or normalized spin-$N/2$ triple.  Its localizer
has one extra positive channel: positive and negative multiplicities are
$N+2$ and $N$.  By contrast, a commuting triple has joint eigenvalues
$z\in\mathbb R^3$, and each nonzero $z$ contributes the balanced pair
$\{+|z|,-|z|\}$ to its localizer.  Since the fuzzy-sphere localizer has a
uniform gap $1/2$ about zero, this inertia mismatch cannot disappear under a
small perturbation.

\section{Counterexample theorem}

\begin{theorem}[uniform obstruction]\label{thm:main}
For every integer $N\ge2$, let $H_{i,N}=T_N(s_i)$, $i=1,2,3$.  Suppose
$Y_{1,N},Y_{2,N},Y_{3,N}\in M_{2^N}(\C)$ are self-adjoint and pairwise
commuting, and suppose each $Y_{i,N}$ is fixed under conjugation by both
$U_{(1\,2)}$ and $U_{(1\,2\,\ldots\,N)}$.  Then
\[
 \max_{1\le i\le3}\|Y_{i,N}-H_{i,N}\|\ge \frac16.
\]
Consequently, the macroscopic Pauli triple admits no asymptotically close
commuting self-adjoint approximants preserving those two symmetries
simultaneously.
\end{theorem}

The proof uses two elementary lemmas.

\begin{lemma}[symmetric-block localizer]\label{lem:spin}
Let $K_N=\Sym^N(\C^2)$ and let $X_i=H_{i,N}|_{K_N}$.  Then
$\dim K_N=N+1$ and the spectrum of $\mathcal L(X)$ is
\[
 \left\{\frac12\ \text{with multiplicity }N+2,\quad
 -\left(\frac12+\frac1N\right)\ \text{with multiplicity }N\right\}.
\]
In particular, $\mathcal L(X)$ is invertible, its distance from singular
self-adjoint matrices is $1/2$, and its inertia is $(N+2,N)$.
\end{lemma}

\begin{proof}
The symmetric power $K_N$ carries the irreducible spin $j=N/2$
representation.  If $J_1,J_2,J_3$ are its generators, then
$X_i=J_i/N$.  We compute the localizer directly, avoiding any index theorem.

Choose the $J_3$-basis $\{|m\rangle:-j\le m\le j\}$, so
\[
 J_3|m\rangle=m|m\rangle,\qquad
 J_-|m\rangle=\sqrt{(j+m)(j-m+1)}\,|m-1\rangle.
\]
Writing $\tau_\pm=(\tau_1\pm i\tau_2)/2$, we have
\[
 N\mathcal L(X)=J_3\otimes\tau_3+J_+\otimes\tau_-+J_-\otimes\tau_+.
\]
The vectors $|j\rangle\otimes|\!\uparrow\rangle$ and
$|-j\rangle\otimes|\!\downarrow\rangle$ are eigenvectors with eigenvalue
$j$.  For each $m=-j+1,\ldots,j$, the span of
$|m\rangle\otimes|\!\downarrow\rangle$ and
$|m-1\rangle\otimes|\!\uparrow\rangle$ is invariant, and the matrix of
$N\mathcal L(X)$ there is
\[
 \begin{pmatrix}
 -m & c_m\\ c_m&m-1
 \end{pmatrix},\qquad
 c_m^2=(j+m)(j-m+1)=j(j+1)-m(m-1).
\]
Its trace is $-1$ and its determinant is $-j(j+1)$, so its eigenvalues are
$j$ and $-(j+1)$.  There are $2j=N$ such two-dimensional blocks, plus the two
extreme eigenvectors.  Dividing by $N=2j$ gives the stated spectrum and
multiplicities.
\end{proof}

\begin{lemma}[commuting triples have balanced inertia]\label{lem:balanced}
If $Z_1,Z_2,Z_3$ are commuting self-adjoint operators on an $n$-dimensional
complex Hilbert space and $\mathcal L(Z)$ is invertible, then its inertia is
$(n,n)$.
\end{lemma}

\begin{proof}
Simultaneously diagonalize the $Z_i$.  For every joint eigenvalue
$z=(z_1,z_2,z_3)\in\mathbb R^3$, the corresponding two-dimensional block of
$\mathcal L(Z)$ is $z_1\tau_1+z_2\tau_2+z_3\tau_3$, whose eigenvalues are
$\pm|z|$.  Invertibility excludes $z=0$, so every joint eigenspace contributes
equally to the positive and negative counts.
\end{proof}

\begin{proof}[Proof of Theorem \ref{thm:main}]
The transposition and long cycle generate $S_N$.  Fixedness under conjugation
by their unitaries therefore implies $[Y_{i,N},U_\sigma]=0$ for every
$\sigma\in S_N$.  The group average
\[
 P_N=\frac1{N!}\sum_{\sigma\in S_N}U_\sigma
\]
is the orthogonal projection onto $K_N=\Sym^N(\C^2)$, hence every $Y_{i,N}$
commutes with $P_N$.  The restrictions
$Z_i=Y_{i,N}|_{K_N}$ are consequently self-adjoint and pairwise commuting,
and
\[
 \|Z_i-X_i\|\le \|Y_{i,N}-H_{i,N}\|.
\]

Assume for contradiction that the maximum on the right is less than $1/6$.
Since $\|\tau_i\|=1$,
\[
 \|\mathcal L(Z)-\mathcal L(X)\|
 \le\sum_{i=1}^3\|Z_i-X_i\|<\frac12.
\]
Lemma \ref{lem:spin} says that $\mathcal L(X)$ has spectral gap $1/2$ about
zero.  Thus the straight-line path from $\mathcal L(X)$ to $\mathcal L(Z)$
remains invertible, so inertia is constant along it.  Hence
$\mathcal L(Z)$ would have inertia $(N+2,N)$.  But the $Z_i$ commute on the
$(N+1)$-dimensional space $K_N$, so Lemma \ref{lem:balanced} gives inertia
$(N+1,N+1)$, a contradiction.
\end{proof}

\section{Checks, limitations, and novelty status}

\paragraph{Independent calculation.}
The included script \texttt{code/verify\_localizer.py} constructs the
spin matrices and checks the exact spectrum, Casimir relation, commutators,
and gap numerically for $1\le N\le24$.  Errors were below
$3\cdot10^{-16}$.  This is only a consistency check; the proof above is exact.

\paragraph{What is and is not disproved.}
The theorem answers negatively the simultaneous preservation of the two
specific linear multiplicative symmetries which the source identifies as
generators of all tensor permutations.  It therefore also rules out
simultaneously preserving all $S_N$ symmetries.  It does not rule out
preservation of one fixed permutation symmetry at a time.  Nor does it decide
the source's transpose, conjugation, real, symplectic, phase, or signed
symmetry patterns.  Ordinary Ogata approximation remains valid: the
counterexample shows only that its approximants must break enough
tensor-permutation symmetry to couple the obstructed sector to other sectors.

\paragraph{Bounded novelty audit.}
The run indexes contained no duplicate.  On 11 August 2026, OpenAlex reported
zero citing works for arXiv:2412.20795.  Exact-phrase searches for the source
question and close searches combining Ogata approximation, tensor-permutation
invariance, symmetric powers, spin matrices, Bott index, and localizers found
no direct answer.  We inspected the source, Ogata's theorem
\cite{Ogata2013}, Herrera's constructive qubit treatment \cite{Herrera2022},
and standard Bott/localizer sources \cite{HastingsLoring,LoringGuide}.
Those sources supply context for the obstruction but do not state this
application to simultaneous permutation symmetry.  The novelty assessment is
therefore moderate-to-high under a bounded search, not an exhaustive priority
claim.

\paragraph{Human-review recommendation.}
Prioritize review.  Verify (i) that simultaneous preservation of the two
generators is within the intended scope of Example 2.5, (ii) the reduction to
$\Sym^N(\C^2)$, and (iii) the two localizer inertia calculations.  No
conditional lemma remains.

\begin{thebibliography}{9}
\bibitem{Herrera2024}
D. Herrera,
\emph{A Projection Characterization and Symmetry Bootstrap for Elements of a
von Neumann Algebra that are Nearby Commuting Elements},
arXiv:2412.20795v1 (2024), Example 2.5.

\bibitem{Ogata2013}
Y. Ogata,
\emph{Approximating macroscopic observables in quantum spin systems with
commuting matrices}, J. Funct. Anal. \textbf{264} (2013), 2005--2033;
arXiv:1111.5933.

\bibitem{Herrera2022}
D. Herrera,
\emph{Constructing Nearby Commuting Matrices for Reducible Representations of
$su(2)$ with an Application to Ogata's Theorem}, arXiv:2212.06012v2 (2023).

\bibitem{HastingsLoring}
M. B. Hastings and T. A. Loring,
\emph{Almost Commuting Matrices, Localized Wannier Functions, and the Quantum
Hall Effect}, J. Math. Phys. \textbf{51} (2010), 015214;
arXiv:0910.5490.

\bibitem{LoringGuide}
T. A. Loring,
\emph{A Guide to the Bott Index and Localizer Index}, arXiv:1907.11791 (2019).
\end{thebibliography}

\end{document}
